
\paragraph{模式 5、6 示例}
\subsubsection*{Example for Diode Modes 5 and 6}

假设天线面积按单层模式表面面积（antenna mode 1）计算，M1 的 diode ratio vector 为 \texttt{\{0 0 1 0\}}（默认），M1 的 \texttt{layerMaxRatio} 为 400，同一 M1 网络连接二极管 A（保护值 0.5）、B（1.0）、C（1.5）；未指定 \texttt{v4} 时按 0 处理。

\begin{lstlisting}
icc2_shell> define_antenna_rule -mode 1 -diode_mode diode_mode \
   -metal_ratio 400 -cut_ratio 20
\end{lstlisting}

因使用全局 \texttt{layerMaxRatio} 与默认向量，无需层相关规则。模式 5/6 的最大天线比为
\texttt{metal\_area / (gate\_area + equi\_gate\_area)}，其中
\texttt{equi\_gate\_area = (diode\_protection + v1) * (v2 + v3)}。
\begin{itemize}
  \item 模式 5：用最大保护值 1.5，得 \texttt{metal\_area/(gate\_area + 1.5)}；
  \item 模式 6：用总保护值 \(0.5+1.0+1.5=3.0\)，得 \texttt{metal\_area/(gate\_area + 3.0)}。
\end{itemize}

\paragraph{模式 7、8 示例}
\subsubsection*{Example for Diode Modes 7 and 8}

假设单层表面面积模式，M1 向量为 \texttt{\{0.7 0.0 150 800\}}，\texttt{layerMaxRatio}=400，二极管保护值同前。

\begin{lstlisting}
icc2_shell> define_antenna_rule -mode 1 -diode_mode diode_mode \
   -metal_ratio 400 -cut_ratio 20
icc2_shell> define_antenna_layer_rule -mode 1 -layer "M1" \
   -ratio 400 -diode_ratio {0.7 0.0 150 800}
\end{lstlisting}

模式 7/8 计算为 \texttt{(metal\_area - equi\_metal\_area) / gate\_area}：
\begin{itemize}
  \item 模式 7：最大保护 1.5，得 \texttt{(metal\_area - 1025) / gate\_area}；
  \item 模式 8：总保护 3.0，得 \texttt{(metal\_area - 1250) / gate\_area}。
\end{itemize}

\paragraph{模式 2、3、4 的网络最大天线比}
对表~\ref{tab:max-ratio-per-diode} 中二极管 A/B/C（向量 \texttt{\{0.7 0.0 200 2000\}}，\texttt{layerMaxRatio}=400）的计算结果为：A$\to$400，B$\to$2200，C$\to$2300。则：
\begin{itemize}
  \item 模式 2：取最大 $=2300$；
  \item 模式 3：求和 $=400+2200+2300=4900$；
  \item 模式 4：用保护值之和代入公式，\((0.5+1.0+1.5)*200+2000=2600\)。
\end{itemize}

定义规则示例：
\begin{lstlisting}
icc2_shell> define_antenna_rule -mode 1 -diode_mode diode_mode \
   -metal_ratio 400 -cut_ratio 20
icc2_shell> define_antenna_layer_rule -mode 1 -layer "M1" \
   -ratio 400 -diode_ratio {0.7 0.0 200 2000}
\end{lstlisting}

\paragraph{模式 14（多栅氧厚度）示例}
\subsubsection*{Example for Diode Mode 14 With Multiple Gate Oxide Thicknesses}

假设单层表面面积模式，设计含两种栅氧厚度。gate class 0 的 M1 向量为
\texttt{\{1 0 1e9 1e9 0 1 2 285 1 285\}}，gate class 1 为
\texttt{\{1 0 1e9 1e9 1 0 2 165 1 28\}}，\texttt{layerMaxRatio}=285。
模式 14 每个 gate class 需 10 个值；即使仅用两类，定义时仍须给出全部 40 个值。未用 class 填很大比率以禁用检查。

\begin{lstlisting}
icc2_shell> define_antenna_rule -mode 1 -diode_mode 14 \
   -metal_ratio 285 -cut_ratio 10
icc2_shell> define_antenna_layer_rule -mode 1 -layer "M1" \
   -ratio 285 -diode_ratio {1 0 1e9 1e9 0 1 2 285 1 285 \
                            1 0 1e9 1e9 1 0 2 165 1 28 \
                            1 0 1e9 1e9 0 1 2 1e9 1 1e9 \
                            1 0 1e9 1e9 0 1 2 1e9 1 1e9 }
\end{lstlisting}

向量各分量含义：
\begin{itemize}
  \item V0：公式 1 使用的栅极面积范围（0=当前 class；1=全部 class 总和）；
  \item V1：连二极管时公式 1 的 diode protection 乘数；
  \item V2：连二极管时公式 1 的偏移；
  \item V3：不连二极管时公式 1 的比率要求；
  \item V4：公式 2 的栅极面积范围（同 V0）；
  \item V5--V7：连二极管时公式 2 的栅极乘数、保护乘数与比率要求；
  \item V8--V9：不连二极管时公式 2 的栅极乘数与天线比要求。
\end{itemize}

模式 14 公式（表~\ref{tab:antenna-ratio-calc} 续）：若连二极管，须同时满足
\texttt{(metal/via area)/gate-area <= (dp*v1+v2)} 与
\texttt{(metal/via area)/(v5*gate-area+v6*dp) <= v7}；
若不连二极管，须同时满足
\texttt{(metal/via area)/gate-area <= v3} 与
\texttt{(metal/via area)/(v8*gate-area) <= v9}。

\subsubsection{计算 Pin 的天线比}
\subsubsection*{Calculating the Antenna Ratio for a Pin}

工具支持多种天线比（antenna-area/gate-area）计算方式。通过指定下列信息控制：
\begin{itemize}
  \item 如何计算天线面积（antenna area mode）；
  \item 计算时考虑哪些金属段（antenna recognition mode）。
\end{itemize}
用 \cmd{define\_antenna\_rule} 与 \cmd{define\_antenna\_layer\_rule} 的 \opt{-mode} 指定（二者的必选选项），见「设置天线模式」。

\subsection{设置天线模式}
\subsubsection*{Setting the Antenna Mode}

天线模式通过以下两方面控制天线比计算：
\begin{itemize}
  \item \textbf{天线面积计算方式（antenna area mode）}
  \begin{itemize}
    \item 表面面积：\(W \times L\)；
    \item 侧壁面积：\((W+L)\times 2\times\text{thickness}\)。
\begin{noteBox}
使用侧壁面积时，须在 technology file 各 Layer 段中用 \texttt{unitMinThickness}、\texttt{unitNomThickness}、\texttt{unitMaxThickness} 定义金属厚度。
\end{noteBox}
  \end{itemize}
  \item \textbf{天线识别模式（antenna recognition mode）}
  \begin{itemize}
    \item \textbf{Single-layer mode}：仅考虑当前层金属段，忽略下层。可布线性最好。\\
          \texttt{antenna\_ratio = 该层相连金属面积 / 总栅极面积}
    \item \textbf{Accumulated-ratio mode}：考虑当前层及与输入 pin 相连的下层段。\\
          \texttt{antenna\_ratio = 当前层与下层各 single-layer 比率之和}
    \item \textbf{Accumulated-area mode}：考虑当前层与下层金属面积。\\
          \texttt{antenna\_ratio = 当前及下层相连金属面积 / 总栅极面积}
  \end{itemize}
\end{itemize}

用 \opt{-mode} 指定天线模式（必选）。表~\ref{tab:antenna-mode} 列出取值。

\begin{table}[htbp]
\centering
\caption{天线模式设置（原书 Table 32）}
\label{tab:antenna-mode}
\begin{tabular}{@{}ccc@{}}
\toprule
天线模式值 & 面积计算方式 & 天线识别模式 \\
\midrule
1 & Surface & Single-layer \\
2 & Surface & Accumulated-ratio \\
3 & Surface & Accumulated-area \\
4 & Sidewall & Single-layer \\
5 & Sidewall & Accumulated-ratio \\
6 & Sidewall & Accumulated-area \\
\bottomrule
\end{tabular}
\end{table}

\subsubsection{天线识别模式示例}
\subsubsection*{Antenna Recognition Mode Example}

图~\ref{fig:antenna-recog-layout} 给出含侧视的布局示例；表~\ref{tab:antenna-recog-ratios} 列出各识别模式下的天线比。

\begin{figure}[htbp]
\centering
\fbox{\parbox{0.85\textwidth}{\centering\vspace{1.5cm}
\textit{（原书 Figure 89：Layout Example for Antenna Recognition Modes）}\\[0.3em]
\small 请参阅原版 PDF 插图。
\vspace{1.5cm}}}
\caption{天线识别模式布局示例}
\label{fig:antenna-recog-layout}
\end{figure}

\begin{table}[htbp]
\centering
\caption{天线识别模式与比率（原书 Table 33）}
\label{tab:antenna-recog-ratios}
\begin{tabular}{@{}lp{0.62\textwidth}@{}}
\toprule
考虑的段 & 天线比 \\
\midrule
Single-layer &
M1：Gate1$=$M1b/Gate1；Gate2$=$M1c/Gate2；Gate3$=$M1d/Gate3。\\
M2：Gate1$=$M2b/Gate1；Gate2$=$M2d/Gate2；Gate3$=$M2e/Gate3。\\
M3：Gate1,2$=$(M3b+M3c)/(Gate1+Gate2)；Gate3$=$M3d/Gate3。\\
M4：Gate1,2,3$=$(M4a+M4b)/(Gate1+Gate2+Gate3)。 \\
\midrule
Accumulated-ratio &
M1 同 single-layer。\\
M2：Gate1$=$M1b/Gate1$+$M2b/Gate1（Gate2/3 类似累加）。\\
M3/M4：在下层比率上继续累加当前层 single-layer 比率。 \\
\midrule
Accumulated-area &
M1 同 single-layer。\\
M2：Gate1$=$(M1b+M2b)/Gate1 等。\\
M3：Gate1,2 为多层金属面积之和 / (Gate1+Gate2)；Gate3 类似。\\
M4：全部金属面积 / (Gate1+Gate2+Gate3)。 \\
\bottomrule
\end{tabular}
\end{table}

\subsection{指定天线属性}
\subsubsection*{Specifying Antenna Properties}

标准单元与 hard macro 的天线属性一般在参考库的 frame view 中定义。可为未在库中定义属性的单元设置默认值：
\begin{itemize}
  \item \appopt{route.detail.default\_diode\_protection}：标准单元输出 pin 缺省 diode protection；
  \item \appopt{route.detail.default\_gate\_size}：标准单元输入 pin 缺省 gate size；
  \item \appopt{route.detail.default\_port\_external\_antenna\_area}：顶层 port 缺省天线面积；
  \item \appopt{route.detail.default\_port\_external\_gate\_size}：顶层 port 缺省 gate size；
  \item \appopt{route.detail.macro\_pin\_antenna\_mode}：macro pin 的天线处理方式；
  \item \appopt{route.detail.port\_antenna\_mode}：顶层 port 的天线处理方式。
\end{itemize}

层次化流程中若为 block 创建 abstraction，须用 \cmd{derive\_hier\_antenna\_property} 从 block 提取天线信息供上一层次使用。

\begin{noteBox}
若 macro 并非所有层都定义了层次化天线属性，Zroute 视为数据不完整，跳过天线分析并给出 ZRT-311 警告。参见 SolvNet 文章 027178。
\end{noteBox}

\begin{seeAlsoBox}
在 hard macro 上标注天线属性（Annotating Antenna Properties on Hard Macro Cells）。
\end{seeAlsoBox}

\subsection{分析并修复天线违例}
\subsubsection*{Analyzing and Fixing Antenna Violations}

若设计库含天线规则，Zroute 在详细布线期间自动分析并修复天线违例。要禁用，将 \appopt{route.detail.antenna} 设为 \cmd{false}。天线规则与其他设计规则一样在详细布线中并发检查与修正，以减少迭代、缩短总运行时间。

默认行为：
\begin{itemize}
  \item 对全部时钟与信号网检查并修复；可用 \appopt{route.detail.skip\_antenna\_fixing\_for\_nets} 禁止对特定网修复（仍分析并报告）；
  \item 不对电源/地网检查或修复；将 \appopt{route.detail.check\_antenna\_on\_pg} 设为 \cmd{true} 可启用；
  \item 在第二次迭代（初始布线完成且基本 DRC 已修复后）开始修复；可用 \appopt{route.detail.antenna\_on\_iteration} 更改；
  \item 通过 \textbf{layer hopping} 修复：把大金属多边形拆到更高层，或（对最高层有限保护情形）下移到更低层。
\end{itemize}

Layer hopping 类型：
\begin{itemize}
  \item 用更高层金属段打断天线：对 metal-N 违例，在靠近栅极处插入小段 metal-(N+1)，降低剩余 metal-N 比率。不适用于最高层且输出保护有限的情况。
  \item 下移到更低层：主要用于最高层违例；可能增加 RC 与时序延迟。
\end{itemize}

也可插入二极管修复。将 \appopt{route.detail.insert\_diodes\_during\_routing} 设为 \cmd{true} 启用。要强制仅用二极管，将 \appopt{route.detail.hop\_layers\_to\_fix\_antenna} 设为 \cmd{false}。若两者皆启用，默认先尝试 layer hopping；要将偏好改为二极管，设 \appopt{route.detail.antenna\_fixing\_preference} 为 \texttt{use\_diodes}。

删除修复天线所不需要的冗余二极管：\appopt{route.detail.delete\_redundant\_diodes\_during\_routing}=\cmd{true}。Spare 二极管不删除；详细布线插入的 port 保护二极管也不删除。

每轮详细布线结束会报告天线违例，例如：
\begin{lstlisting}
DRC-SUMMARY:
   @@@@@@@ TOTAL VIOLATIONS =      506
    @@@@ Total number of instance ports with antenna violations = 1107
\end{lstlisting}

检查天线违例：
\begin{lstlisting}
icc2_shell> check_routes -antenna true
\end{lstlisting}

\subsection{详细布线期间插入二极管}
\subsubsection*{Inserting Diodes During Detail Routing}

通过在靠近被保护栅极的网络上插入反向偏置二极管，为累积电荷提供放电通路，且正常工作时不应导通电流。

启用：\appopt{route.detail.insert\_diodes\_during\_routing}=\cmd{true}。

控制选项：
\begin{itemize}
  \item \appopt{route.detail.antenna\_fixing\_preference}=\texttt{use\_diodes}：优先二极管而非 hopping；
  \item \appopt{route.detail.hop\_layers\_to\_fix\_antenna}=\cmd{false}：强制二极管修复；
  \item 默认可新增二极管或使用 spare；用 \appopt{route.detail.diode\_preference} 设为 \texttt{new}/\texttt{spare}/\texttt{none}；
  \item \appopt{route.detail.diode\_insertion\_mode}=\texttt{new}/\texttt{spare}/\texttt{new\_and\_spare}：强制方法；
\begin{noteBox}
要利用 spare 二极管，须在标准单元布局前/后、天线违例可能出现区域布线之前加入 spare。
\end{noteBox}
  \item \appopt{route.detail.diode\_libcell\_names}：限制可用二极管（仅写单元名，勿含库名）；
  \item \appopt{route.detail.reuse\_filler\_locations\_for\_diodes}=\cmd{false}：禁止复用 filler 位置。
\end{itemize}

层次与 voltage area：
\begin{itemize}
  \item Pin 违例：二极管插在与违例 pin 相同逻辑层次；
  \item 顶层 port 违例：默认顶层；若 port 属 voltage area，可将 \appopt{route.detail.use\_lower\_hierarchy\_for\_port\_diodes} 设为 \cmd{true}。
\end{itemize}

###{详细布线后插入二极管}
\subsubsection*{Inserting Diodes After Detail Routing}

详细布线后可用 \cmd{create\_diodes} 显式指定违例与约束。根据约束与 \appopt{route.detail.diode\_insertion\_mode}，插入新二极管或复用 spare。

须用列表指定：port、单元实例、二极管参考单元、数量、连接最高层、相对 pin 的最大距离。距离内无法插入/复用则跳过。

\begin{lstlisting}
{port_name instance_name diode_reference number_of_diodes
   max_routing_layer max_routing_distance}
\end{lstlisting}

\begin{noteBox}
将单元实例指定为顶层 block 名，可为顶层 port 插入二极管。
\end{noteBox}

相关应用选项：\appopt{route.detail.diode\_insertion\_mode}、\appopt{route.detail.reuse\_filler\_locations\_for\_diodes}、\appopt{route.detail.use\_lower\_hierarchy\_for\_port\_diodes}。

示例（CI 的 A port，2 个 Adiode，最高 M5，2.5\,µm）：
\begin{lstlisting}
icc2_shell> create_diodes {{A CI Adiode 2 M5 2.5}}
\end{lstlisting}

% ============================================================
\section{插入冗余 Via}
\subsection*{Inserting Redundant Vias}

冗余 via 插入是 Zroute 在整个布线流程中支持的重要 DFM 功能。各布线阶段并发优化 via 数量与线长。结果以冗余 via 转换率（单 via 转为冗余 via 的百分比）衡量；亦应关注未优化单 via 数量——越少通常对 DFM 越有利。

相关主题：
\begin{itemize}
  \item 在时钟网上插入冗余 Via
  \item 在信号网上插入冗余 Via
\end{itemize}

\subsection{在时钟网上插入冗余 Via}
\subsubsection*{Inserting Redundant Vias on Clock Nets}

Zroute 可在时钟布线期间或之后插入冗余 via。布线后插入见「布线后冗余 Via 插入」。

时钟布线期间步骤：
\begin{enumerate}
  \item 用 \cmd{create\_routing\_rule} 在非默认布线规则中指定冗余 via（见「指定非默认 Via」）。选择时须兼顾 DFM 与可布线性。NDR 还可定义更严线宽/间距与 tapering 距离。
  \item 用 \cmd{set\_clock\_routing\_rules} 将 NDR 赋给时钟网：
\begin{lstlisting}
icc2_shell> set_clock_routing_rules -rules clock_via_rule
\end{lstlisting}
  \item 运行 CTS：
\begin{lstlisting}
icc2_shell> synthesize_clock_trees
\end{lstlisting}
  \item 布线时钟网：
\begin{lstlisting}
icc2_shell> route_group -all_clock_nets \
   -reuse_existing_global_route true
\end{lstlisting}
        全局布线预留冗余 via 空间，详细布线插入冗余 via。
\end{enumerate}

\subsection{在信号网上插入冗余 Via}
\subsubsection*{Inserting Redundant Vias on Signal Nets}

可采用：
\begin{itemize}
  \item 布线后冗余 via 插入；
  \item 并发 soft-rule 冗余 via 插入；
  \item 接近 100\% 冗余 via 插入。
\end{itemize}

一般先做布线后插入；若转换率 $\ge$80\% 可再尝试并发 soft-rule；若 $\ge$90\% 可再尝试接近 100\% 插入。

\begin{noteBox}
转换率升高会使设计规则更难收敛，并可能降低信号完整性；接近 100\% 仅用于确需极高转换率的 block，且可能需调整 floorplan utilization 以留出空间。
\end{noteBox}

相关主题：查看默认 Via 映射表；定义定制 Via 映射表；布线后插入；并发 soft-rule；接近 100\%；时序保持；报告转换率。

\subsubsection{查看默认 Via 映射表}
\subsubsection*{Viewing the Default Via Mapping Table}

默认 Zroute 从 technology file 读取默认 contact code 并生成优化映射表。查看：
\begin{lstlisting}
icc2_shell> add_redundant_vias -list_only true
\end{lstlisting}
多数情况下定制映射表优于默认表。若已用 \cmd{add\_via\_mapping} 创建定制表，该命令显示定制表。

示例（原书 Example 25）：
\begin{lstlisting}
icc2_shell> add_redundant_vias -list_only
...
Redundant via optimization will attempt to replace the following vias:
    VIA12SQ_C    -> VIA12SQ_C_2x1    VIA12SQ_C_2x1(r) VIA12SQ_C_1x2 ...
    VIA12SQ_C(r) -> VIA12SQ_C_2x1    VIA12SQ_C_2x1(r) VIA12SQ_C_1x2 ...
...
\end{lstlisting}

\subsubsection{定义定制 Via 映射表}
\subsubsection*{Defining a Customized Via Mapping Table}

使用 \cmd{add\_via\_mapping}。至少用 \opt{-from}/\opt{-to} 指定源 via 与替换 via（须为 technology file 中的 via 或 \cmd{create\_via\_def} 创建的设计专用 via）。\opt{-from} 须为 simple via 或 via array；\opt{-to} 可为 simple、simple array 或 custom via。

细化选项：
\begin{itemize}
  \item \opt{-weight}：优先级 1--10（默认 1）；高权重优先使用；
  \item \opt{-transform}：\texttt{all}（默认可旋转/翻转）、\texttt{rotate}、\texttt{flip}、\texttt{none}。
\end{itemize}

映射保存在设计库 via 映射表中。覆盖已有映射须 \opt{-force}。

\begin{lstlisting}
icc2_shell> add_via_mapping \
   -from {VIA12 1x1} -to {VIA12T 2x1} -weight 5
icc2_shell> add_via_mapping \
   -from {VIA12 1x1} -to {VIA12 2x1} -weight 1
\end{lstlisting}

用 \cmd{report\_via\_mapping} 查看（\opt{-from}/\opt{-to} 可过滤）。用 \cmd{remove\_via\_mappings} 移除（\opt{-all}/\opt{-from}/\opt{-to}）。

\subsubsection{使用 Via 映射表子集}
\subsubsection*{Using a Subset of the Via Mapping Table for Redundant Via Insertion}

通过应用选项按层选择权重组：
\begin{itemize}
  \item \appopt{route.common.redundant\_via\_include\_weight\_group\_by\_layer\_name}：仅用指定权重组；未列出的层不做插入；
  \item \appopt{route.common.redundant\_via\_exclude\_weight\_group\_by\_layer\_name}：排除指定权重组；
  \item 两者同时设置时，在 include 集合中再排除 exclude 集合。
\end{itemize}

\subsubsection{布线后冗余 Via 插入}
\subsubsection*{Postroute Redundant Via Insertion}

使用 \cmd{add\_redundant\_vias}：可将单 cut 换成多 cut 阵列、换成其他 contact code 的单 cut，或换成不同多 cut 阵列。插入时详细布线检查邻近 partition 的设计规则以减少违例。

默认对所有网插入；\opt{-nets} 可限制网络。完成后重新检查并修复 DRC。

若转换率不足，可用 \opt{-effort high}（约再提升 3--5\%，通过移动 via 腾空间）；在 45\,nm 及以下可能不利于光刻友好图案，此时宜改用并发 soft-rule。也可将 \appopt{route.detail.optimize\_wire\_via\_effort\_level} 设为 \texttt{high} 以减少单 via、缩短线长。

初始布线后插入完成后，设 \appopt{route.common.post\_detail\_route\_redundant\_via\_insertion}，以便后续详细/ECO 布线后自动插入，维持转换率。

\subsubsection{并发 Soft-Rule 冗余 Via 插入}
\subsubsection*{Concurrent Soft-Rule-Based Redundant Via Insertion}

在布线期间为冗余 via 预留空间以提高转换率；可用于初始与 ECO 布线。实际插入仍须事后运行 \cmd{add\_redundant\_vias}。

\begin{noteBox}
预留空间会增加布线运行时间；仅当布线后方法转换率已 $\ge$80\% 且需进一步提高时使用。
\end{noteBox}

步骤：
\begin{enumerate}
  \item （可选）定义 via 映射表；
  \item 启用并发 soft-rule：
\begin{lstlisting}
icc2_shell> set_app_options \
   -name route.common.concurrent_redundant_via_mode \
   -value reserve_space
\end{lstlisting}
        用 \appopt{route.common.concurrent\_redundant\_via\_effort\_level} 控制努力（默认 low；high 也影响详细布线，可能影响 DRC 收敛）。若在 \cmd{place\_opt} 前启用，拥塞估计会考虑冗余 via。\\
        ECO：设 \appopt{route.common.eco\_route\_concurrent\_redundant\_via\_mode}=\texttt{reserve\_space}，并用对应 effort 选项。
\begin{noteBox}
ECO 期间使用可能影响时序与 DRC 收敛；一般仅当初始布线用了接近 100\% 插入时才在 ECO 用 soft-rule 维持转换率。
\end{noteBox}
  \item 布线 block（预留空间并修复硬规则）；
  \item 运行布线后 \cmd{add\_redundant\_vias}。
\end{enumerate}

\subsubsection{接近 100\% 冗余 Via 插入}
\subsubsection*{Near 100 Percent Redundant Via Insertion}

将冗余 via 作为硬设计规则处理。仅支持初始布线，不支持 ECO；初始用此法后，ECO 应改用 soft-rule 以保持转换率。

\begin{noteBox}
对拥塞 block 可能大幅增加运行时间；仅当 soft-rule 已达 $\ge$90\% 且仍需提高时使用。
\end{noteBox}

\begin{lstlisting}
icc2_shell> set_app_options \
   -name route.common.concurrent_redundant_via_mode \
   -value insert_at_high_cost
\end{lstlisting}

布线期间插入冗余 via 并修复硬规则；优先级通常与其他硬规则相同，若详细布线 DRC 不收敛会自动放宽冗余 via 约束。

\subsubsection{冗余 Via 插入时保持时序}
\subsubsection*{Preserving Timing During Redundant Via Insertion}

插入冗余 via 会改变时序：短网因电容增加变慢，长网因电阻下降变快。时序保持模式通过对关键网禁止插入来避免影响时序。

在 \cmd{add\_redundant\_vias} 上使用：
\begin{itemize}
  \item \opt{-timing\_preserve\_setup\_slack\_threshold}
  \item \opt{-timing\_preserve\_hold\_slack\_threshold}
  \item \opt{-timing\_preserve\_nets}
\end{itemize}

应在流程末尾、常规插入已收敛转换率与时序 QoR 之后使用；过早使用可能因关键网过多而严重降低转换率。

\subsection{报告冗余 Via 率}
\subsubsection*{Reporting Redundant Via Rates}

插入后 Zroute 生成报告，包含：
\begin{itemize}
  \item 非默认 via 的总优化转换率；
  \item 各层优化转换率（含 double via 与 DFM-friendly bar via），并分别相对总 via 数与可优化 routed via 数给出；
\begin{noteBox}
使用 bar via 时优化转换率参考价值有限。
\end{noteBox}
  \item 各层按 weight 的优化 via 分布（高于某权重的转换率须对各权重求和）；无权重记为 ``Un-optimized''；
  \item block 总 double via 转换率。
\end{itemize}

\begin{noteBox}
\cmd{report\_design} 的冗余 via 率与 Zroute/\cmd{check\_routes} 不同：前者含 PG 与信号 via，后者仅信号 via，且基于冗余 via 映射。
\end{noteBox}

示例（原书 Example 27）：
\begin{lstlisting}
Total optimized via conversion rate = 96.94% (1401030 / 1445268 vias)
Layer V01        = 41.89% (490617 / 1171301 vias)
    Weight 10    = 9.64% (112869 vias)
    Weight 5     = 32.25% (377689 vias)
    ...
Total double via conversion rate = 46.69% (2158006 / 4622189 vias)
\end{lstlisting}

% ============================================================
\section{优化线长与 Via 数量}
\subsection*{Optimizing Wire Length and Via Count}

详细布线时，Zroute 在存在 DRC 违例的区域优化线长与 via 数量，但不优化无违例区域。

为提高良率，在详细布线与冗余 via 插入之后，可用 \cmd{optimize\_routes} 独立优化线长与 via 数量。

默认按总体代价选择重布网络；对所选网决定重布全部形状或部分形状。用 \opt{-nets} 指定网络；可用 \opt{-reroute\_all\_shapes\_in\_nets} 强制重布该网全部形状。

默认最多 40 次详细布线迭代以修复优化后违例；可用 \opt{-max\_detail\_route\_iterations} 控制。

% ============================================================
\section{减小临界面积}
\subsection*{Reducing Critical Areas}

临界面积是这样一块区域：若随机颗粒缺陷中心落在该处，会导致电路失效、降低良率。导电缺陷造成短路，非导电缺陷造成开路。

相关主题：
\begin{itemize}
  \item 通过 wire spreading 减小短路临界面积；
  \item 通过 wire widening 减小开路临界面积。
\end{itemize}

\subsection{执行 Wire Spreading}
\subsubsection*{Performing Wire Spreading}

详细布线与冗余 via 插入后，可用 \cmd{spread\_wires} 增大平均线间距，减小短路临界面积。默认对同层信号线：
\begin{itemize}
  \item 在优选方向疏导半个 pitch（用 \opt{-pitch} 以层 pitch 倍数修改，如 \texttt{-pitch 1.5}）；
  \item 最小 jog 长度为两倍层 pitch（\opt{-min\_jog\_length}，层 pitch 整数倍）；
  \item 最小 jog 间距为最小层间距加半个层 pitch（\opt{-min\_jog\_spacing\_by\_layer\_name}，按层以微米指定）。
\end{itemize}

图~\ref{fig:wire-spreading} 示出 jog 长度与间距的用法。

\begin{figure}[htbp]
\centering
\fbox{\parbox{0.85\textwidth}{\centering\vspace{1.5cm}
\textit{（原书 Figure 90：Wire Spreading Results）}\\[0.3em]
\small 请参阅原版 PDF 插图。
\vspace{1.5cm}}}
\caption{Wire Spreading 结果}
\label{fig:wire-spreading}
\end{figure}

\begin{lstlisting}
icc2_shell> spread_wires -min_jog_length 3 \
   -min_jog_spacing_by_layer_name {{M1 0.07} {M2 0.08}}
\end{lstlisting}

疏导后自动进行详细布线迭代修复 DRC。时序保持模式选项：
\opt{-timing\_preserve\_setup\_slack\_threshold}、
\opt{-timing\_preserve\_hold\_slack\_threshold}、
\opt{-timing\_preserve\_nets}。
仅对 slack 不小于阈值的网、\opt{-timing\_preserve\_nets} 所列网，以及同层两 pitch 内的相邻网执行疏导。

\subsection{执行 Wire Widening}
\subsubsection*{Performing Wire Widening}

在详细布线、冗余 via 与 wire spreading 之后，用 \cmd{widen\_wires} 增大平均线宽以减小开路临界面积。默认将所有线加宽到原宽 1.5 倍。可用 \opt{-widen\_widths\_by\_layer\_name} 为每层定义最多五种候选宽度：

\begin{lstlisting}
icc2_shell> widen_wires \
   -widen_widths_by_layer_name {{M1 0.07 0.06} {M2 0.08 0.07}}
\end{lstlisting}

加宽会减小邻线间距，可能削弱 spreading 带来的短路临界面积改善。用 \opt{-spreading\_widening\_relative\_weight} 权衡：默认等权；0.0--0.5 偏向 widening（开路）；0.5--1.0 偏向 spreading（短路）。

加宽后进行详细布线迭代修复 DRC；加宽导线不触发 fat wire 间距规则。时序保持选项同 spreading；不对 slack 小于阈值或指定保护网执行加宽。

% ============================================================
\section{插入金属--绝缘体--金属电容}
\subsection*{Inserting Metal-Insulator-Metal Capacitors}

MiM 电容由两层专用导电层夹绝缘体构成，插在两层常规金属之间，见图~\ref{fig:mim-cross-section}。

\begin{figure}[htbp]
\centering
\fbox{\parbox{0.85\textwidth}{\centering\vspace{1.5cm}
\textit{（原书 Figure 91：Cross-Section View of MiM Capacitor Layers）}\\[0.3em]
\small 请参阅原版 PDF 插图（含 M8/MBOT/MTOP/M9 与 viaMimbottom/viaMimtop 等 maskName）。
\vspace{1.5cm}}}
\caption{MiM 电容层剖面示意}
\label{fig:mim-cross-section}
\end{figure}

MiM 电容通常接在电源与地之间，以在电气噪声下维持供电稳定。尺寸常按电源规划的 power strap 几何定制。

用 \cmd{create\_mim\_capacitor\_array} 插入阵列并连接到电源轨。至少指定 MiM 库单元及阵列 $x$/$y$ 增量：

\begin{lstlisting}
icc2_shell> create_mim_capacitor_array \
            -lib_cell my_lib/mim_ref_cell \
            -x_increment 20 -y_increment 20
\end{lstlisting}

默认行为：
\begin{itemize}
  \item 对整个 block 插入；可用 \opt{-boundary} 限制到单个矩形/直角多边形区域。插入时不考虑标准单元、macro、placement blockage 或 voltage area；
  \item 方向 R0；可用 \opt{-orientation}（作用于阵列全部单元）；
  \item 命名：\texttt{mimcap!library\_cell\_name!number}；可用 \opt{-prefix}。
\end{itemize}

% ============================================================
\section{插入 Filler Cell}
\subsection*{Inserting Filler Cells}

为保证电源网连通，可在标准单元行空隙中插入 filler；亦常用于加入去耦电容以稳定供电。工具支持单高与多高 filler。

\begin{noteBox}
插入前须用 \cmd{check\_legality} 确认 block 已合法化。
\end{noteBox}

两种主要方法：
\begin{itemize}
  \item \textbf{有序 filler 参考列表}（标准流程）：按列表顺序在每个间隙插入第一个能放入的单元。若工艺要求特定 filler 邻接标准单元，可先插入这些必选 filler，再用标准流程填剩余空隙；
  \item \textbf{阈值电压规则}（VT-based）：按间隙两侧单元的 VT 类型用用户规则选 filler，多用于成熟晶圆厂节点。
\end{itemize}

\subsection{标准 Filler Cell 插入}
\subsubsection*{Standard Filler Cell Insertion}

使用 \cmd{create\_stdcell\_fillers} 插入，用 \cmd{remove\_stdcell\_fillers\_with\_violation} 检查并移除违例 filler。插入时通过布局合法化器遵守先进节点规则与物理约束；若 \appopt{place.legalize.enable\_advanced\_legalizer}=\cmd{true}，使用先进合法化算法做 2D 规则与单元交互检查，可缩短运行时间。

步骤：
\begin{enumerate}
  \item 用 \cmd{check\_legality} 确认合法。注意：\cmd{check\_legality} 假定所有 \texttt{design\_type==filler} 的标准单元都可用于填充，而 \cmd{create\_stdcell\_fillers} 仅插入 \opt{-lib\_cells} 所列单元，二者差异可能导致插入后仍有不可填间隙。
  \item （可选）用 \cmd{create\_left\_right\_filler\_cells} 为特定标准单元两侧插入必选 filler（见「用特定 Filler 邻接标准单元」）。
  \item （可选）\cmd{set\_host\_options -max\_cores n} 启用多核（仅当 technology file 含最小阈值电压、OD 或 TPO 规则时生效）。
  \item 用 \cmd{create\_stdcell\_fillers -lib\_cells} 插入 metal filler（宜从大到小）。可用：
\begin{lstlisting}
icc2_shell> get_lib_cells ref_lib/* -filter "design_type==filler"
icc2_shell> set FILLER_CELLS \
   [get_object_name [sort_collection -descending \
      [get_lib_cells ref_lib/* -filter "design_type==filler"] area]]
\end{lstlisting}
  \item \cmd{connect\_pg\_net -automatic} 连接 PG。
  \item \cmd{remove\_stdcell\_fillers\_with\_violation} 移除违例 metal filler（可能需多次，因移除会暴露新违例）。
  \item 再插入 nonmetal filler 并连接 PG。
\end{enumerate}

\begin{seeAlsoBox}
Generic ECO Flow for Timing or Functional Changes；移除 Filler Cell。
\end{seeAlsoBox}

\subsection{控制标准 Filler Cell 插入}
\subsubsection*{Controlling Standard Filler Cell Insertion}

默认填充整个 block 水平标准单元行的全部空隙。可控制：
\begin{itemize}
  \item \textbf{按最低漏电选择}：\opt{-leakage\_vt\_order} 按漏电递减给出 VT 层；\opt{-leakage\_vt\_check} 仅检查既有 filler 是否可换成更低漏电单元（最多报告数由 \appopt{chf.create\_stdcell\_fillers.max\_leakage\_vt\_order\_violations} 控制）。须启用 advanced legalizer，且 technology file 中 VT 层 \texttt{maskType=implant}。
  \item \textbf{填充百分比}：\opt{-utilization}；用 \opt{-type\_utilization} 控制各类（如 ULVT/LVT DCAP）比例（总和 $\le$100），可用 \opt{-fill\_remaining} 填剩余，\opt{-prefer\_type\_ordering} 按组内从左到右优先。
  \item \textbf{区域}：\opt{-bboxes}、\opt{-voltage\_area}。
  \item \textbf{电源网违例}：默认不查（省时）；metal filler 建议 \opt{-rules check\_pnet}。
  \item \textbf{防间隙}：\opt{-smallest\_cell\_size} 为 1/2/3 倍 unit site（advanced legalizer 开启时忽略）。
  \item \textbf{Hard placement blockage}：默认遵守；\opt{-ignore\_hard\_blockages} 可忽略。
  \item \textbf{命名}：\texttt{xofiller!...}；\opt{-prefix}/\opt{-separator}。
  \item \textbf{合法化错误}：默认遇错失败；\opt{-continue\_on\_error} 可强制完成。
\end{itemize}

示例：
\begin{lstlisting}
icc2_shell> create_stdcell_fillers -lib_cells $FILLER_CELLS \
   -type_utilization { {*/DCAP16*ULVT */DCAP8*ULVT ...} 70
                       {*/DCAP16*LVT */DCAP8*LVT ...} 30 } \
   -prefer_type_ordering -fill_remaining
\end{lstlisting}

\subsection{检查 Filler Cell DRC 违例}
\subsubsection*{Checking for Filler Cell DRC Violations}

默认 \cmd{remove\_stdcell\_fillers\_with\_violation} 检查名称含 \texttt{xofiller} 的实例与顶层信号布线对象之间的 DRC，并移除违例者。可控制：
\begin{itemize}
  \item \opt{-name}：识别 filler 的字符串；
  \item \opt{-boundary}：检查区域；
  \item \opt{-check\_between\_fixed\_objects true}：与所有邻接对象检查；
  \item \opt{-shorts\_only true}：仅因与布线形状短路而移除；
  \item 双重图案化：当 \appopt{route.common.color\_based\_dpt\_flow}=\cmd{true} 且 technology file 有相应规则时检查；
  \item \opt{-check\_only true}：只检查不移除，写 \texttt{block\_fillers\_with\_violation.rpt}。
\end{itemize}

\subsection{修复剩余 Mask 设计规则违例}
\subsubsection*{Fixing Remaining Mask Design Rule Violations}

部分 mask 规则因出现概率低、计算重而未在插入时处理。插入后用 \cmd{replace\_fillers\_by\_rules -replacement\_rule} 修复。表~\ref{tab:replace-fillers-rules} 列出支持的规则与关键字（原书 Table 34）。

\begin{table}[htbp]
\centering
\caption{\cmd{replace\_fillers\_by\_rules} 支持的设计规则（原书 Table 34）}
\label{tab:replace-fillers-rules}
\begin{tabular}{@{}llp{0.42\textwidth}@{}}
\toprule
设计规则 & \texttt{-replacement\_rule} & 说明 \\
\midrule
End orientation & \texttt{end\_orientation} & 翻转最左/最右 filler 方向 \\
Half row adjacency & \texttt{half\_row\_adjacency} & 半高 filler 邻接限制 \\
Illegal abutment & \texttt{illegal\_abutment} & 禁止某些 filler 邻接 \\
Maximum horizontal length & \texttt{od\_horiztonal\_distance} & 连续 filler 行水平长度 \\
Maximum horizontal edge length & \texttt{horizontal\_edges\_distance} & 特定层邻接对象水平边长 \\
Maximum stacking for small fillers & \texttt{small\_filler\_stacking} & 小 filler 堆叠高度 \\
Maximum vertical edge length & \texttt{max\_vertical\_constraint} & 堆叠 filler 垂直边长 \\
Tap cell spacing & \texttt{od\_tap\_distance} & Tap OD 与邻接 OD 距离 \\
Random & \texttt{random} & 随机替换为定制 filler \\
\bottomrule
\end{tabular}
\end{table}

\subsubsection{End Orientation 规则}
\subsubsection*{End Orientation Rule}

对连续水平单元序列，修改最左/最右且属于目标列表的 filler 方向。须指定 \opt{-target\_fillers}，并用 \opt{-left\_end}/\opt{-right\_end}（\texttt{inverse}、\texttt{R0\_or\_MX}、\texttt{MY\_or\_R180}）。\opt{-check\_only} 仅报告将改方向的单元。

\subsubsection{Half Row Adjacency 规则}
\subsubsection*{Half Row Adjacency Rule}

按邻接标准单元类型（Inbound：顶/底边界在半高处，需 ICIP；Regular：全高边界，邻接侧不得有 ICIP）替换半高 filler。替换单元与原 filler 同高、同宽、同 VT；库提供带 ICIP 的 p/n 型。图~\ref{fig:half-height-filler}；表~\ref{tab:half-height-replace}。

\begin{figure}[htbp]
\centering
\fbox{\parbox{0.85\textwidth}{\centering\vspace{1.2cm}
\textit{（原书 Figure 92：Half-Height Filler Cell Replacement With Half Row Adjacency Rule）}\\[0.3em]
\small 请参阅原版 PDF 插图。
\vspace{1.2cm}}}
\caption{半高 Filler 替换与半行邻接规则}
\label{fig:half-height-filler}
\end{figure}

\begin{table}[htbp]
\centering
\caption{半高替换 Filler Cell（原书 Table 35）}
\label{tab:half-height-replace}
\begin{tabular}{@{}ll@{}}
\toprule
邻接 inbound 标准单元 & 指定替换 filler 的选项 \\
\midrule
无 & \opt{-p\_none}、\opt{-n\_none} \\
仅左 & \opt{-p\_left}、\opt{-n\_left} \\
仅右 & \opt{-p\_right}、\opt{-n\_right} \\
左与右 & \opt{-p\_left\_right}、\opt{-n\_left\_right} \\
顶或底（仅 2X filler） & \opt{-p\_left\_center\_right}、\opt{-n\_left\_center\_right} \\
\bottomrule
\end{tabular}
\end{table}

每个列表每种尺寸只能含一个单元。

\subsubsection{Illegal Abutment 规则}
\subsubsection*{Illegal Abutment Rule}

\begin{lstlisting}
icc2_shell> replace_fillers_by_rules -replacement_rule illegal_abutment \
   -replace_abutment { {FILL1a CFILL1a} {FILL1b CFILL1b} } \
   -illegal_abutment {FILLx FILLy}
\end{lstlisting}

\opt{-skip\_fixed\_cells} 可忽略 \texttt{physical\_status} 为 fixed/locked 的 filler。图~\ref{fig:illegal-abutment}。

\begin{figure}[htbp]
\centering
\fbox{\parbox{0.85\textwidth}{\centering\vspace{1.2cm}
\textit{（原书 Figure 93：Fixing an Illegal Abutment Violation）}\\[0.3em]
\small 请参阅原版 PDF 插图。
\vspace{1.2cm}}}
\caption{修复非法邻接违例}
\label{fig:illegal-abutment}
\end{figure}

\subsubsection{最大水平边长规则}
\subsubsection*{Maximum Horizontal Edge Length Rule}

\opt{-replacement\_rule horizontal\_edges\_distance}，须指定 \opt{-max\_constraint\_length}、\opt{-layer}，以及 \opt{-refill\_table} 或 \opt{-constraint\_fillers}/\opt{-non\_constraint\_fillers}。

\begin{lstlisting}
icc2_shell> replace_fillers_by_rules \
   -replacement_rule horizontal_edges_distance \
   -max_constraint_length 30 -layer M2 \
   -refill_table { {{VT1_FILL1} {VT1_DECAP4 VT1_DECAP2}}
                   {{VT2_FILL1} {VT2_DECAP4 VT2_DECAP2}} }
\end{lstlisting}

\subsubsection{最大水平长度规则}
\subsubsection*{Maximum Horizontal Length Rule}

\opt{-replacement\_rule od\_horiztonal\_distance}（原文关键字拼写如此）。示例限制连续 filler 水平长度 70\,µm：

\begin{lstlisting}
icc2_shell> replace_fillers_by_rules \
   -replacement_rule od_horiztonal_distance -max_constraint_length 70 \
   -refill_table { {{VT1_FILL1} {VT1_DECAP4 VT1_DECAP2}}
                   {{VT2_FILL1} {VT2_DECAP4 VT2_DECAP2}} }
\end{lstlisting}

图~\ref{fig:max-vedge}（原书 Figure 94）。

\begin{figure}[htbp]
\centering
\fbox{\parbox{0.85\textwidth}{\centering\vspace{1.2cm}
\textit{（原书 Figure 94：相关最大长度违例修复示意）}\\[0.3em]
\small 请参阅原版 PDF 插图。
\vspace{1.2cm}}}
\caption{修复最大水平长度相关违例}
\label{fig:max-vedge}
\end{figure}

\subsubsection{小 Filler 最大堆叠规则}
\subsubsection*{Maximum Stacking for Small Filler Cells Rule}

\opt{-replacement\_rule small\_filler\_stacking}，指定 \opt{-max\_constraint\_length}、\opt{-small\_fillers}、\opt{-replacement\_fillers}。超出高度时用大 filler 替换相邻 2--3 个小 filler；无法替换则重排。保持违例处 VT，非违例处 VT 可能改变且不检查。

\begin{lstlisting}
icc2_shell> replace_fillers_by_rules \
   -replacement_rule small_filler_stacking -max_constraint_length 10 \
   -small_fillers {FILL1 FILL2} -replacement_fillers {FILL3 FILL4}
\end{lstlisting}

图~\ref{fig:max-stack-replace}、\ref{fig:max-stack-reorder}。

\begin{figure}[htbp]
\centering
\fbox{\parbox{0.85\textwidth}{\centering\vspace{1.2cm}
\textit{（原书 Figure 95：Fixing Maximum Stacking Violation by Replacement）}\\[0.3em]
\small 请参阅原版 PDF 插图。
\vspace{1.2cm}}}
\caption{通过替换修复最大堆叠违例}
\label{fig:max-stack-replace}
\end{figure}

\begin{figure}[htbp]
\centering
\fbox{\parbox{0.85\textwidth}{\centering\vspace{1.2cm}
\textit{（原书 Figure 96：Fixing Maximum Stacking Violation by Reordering）}\\[0.3em]
\small 请参阅原版 PDF 插图。
\vspace{1.2cm}}}
\caption{通过重排修复最大堆叠违例}
\label{fig:max-stack-reorder}
\end{figure}

\subsubsection{最大垂直边长规则}
\subsubsection*{Maximum Vertical Edge Length Rule}

\opt{-replacement\_rule max\_vertical\_constraint}。指定 \opt{-max\_constraint\_length}、\opt{-constraint\_fillers}（默认所有标准单元受约束，\opt{-exception\_cells} 排除；\opt{-no\_violation\_along\_exception\_cells}）、以及 \opt{-non\_constraint\_fillers}/\opt{-non\_constraint\_left\_fillers}/\opt{-non\_constraint\_right\_fillers}。

\begin{lstlisting}
icc2_shell> replace_fillers_by_rules \
   -replacement_rule max_vertical_constraint -max_constraint_length 180 \
   -constraint_fillers {DECAP8 DECAP4} -non_constraint_fillers {FILL1} \
   -exception_cells {BUF1 TAP1}
\end{lstlisting}

图~\ref{fig:max-vedge2}。

\begin{figure}[htbp]
\centering
\fbox{\parbox{0.85\textwidth}{\centering\vspace{1.2cm}
\textit{（原书 Figure 97：Fixing a Maximum Vertical Edge Length Violation）}\\[0.3em]
\small 请参阅原版 PDF 插图。
\vspace{1.2cm}}}
\caption{修复最大垂直边长违例}
\label{fig:max-vedge2}
\end{figure}

\subsubsection{Tap Cell Spacing 规则}
\subsubsection*{Tap Cell Spacing Rule}

\opt{-replacement\_rule od\_tap\_distance}。指定 \opt{-tap\_distance\_range}（site 数）、\opt{-tap\_cells}，以及 \opt{-left\_violation\_tap}/\opt{-right\_violation\_tap}/\opt{-both\_violation\_tap}。目标与替换 tap 须同尺寸。对无 OD 邻接单元用 \opt{-adjacent\_non\_od\_cells} 跳过。

\begin{lstlisting}
icc2_shell> replace_fillers_by_rules -replacement_rule od_tap_distance \
   -tap_cells {TAP} -tap_distance_range {3 3} \
   -left_violation_tap {TAPl} -right_violation_tap {TAPr} \
   -both_violation_tap {TAPb}
\end{lstlisting}

图~\ref{fig:tap-spacing-fix}。

\begin{figure}[htbp]
\centering
\fbox{\parbox{0.85\textwidth}{\centering\vspace{1.2cm}
\textit{（原书 Figure 98：Fixing a Tap Cell Spacing Violation）}\\[0.3em]
\small 请参阅原版 PDF 插图。
\vspace{1.2cm}}}
\caption{修复 Tap Cell 间距违例}
\label{fig:tap-spacing-fix}
\end{figure}

\subsection{随机 Filler Cell 替换}
\subsubsection*{Random Filler Cell Replacement}

\begin{lstlisting}
icc2_shell> replace_fillers_by_rules -replacement_rule random \
   -random_replace { {FILL1 {FILL1a FILL1b FILL1c}}
                     {FILL2 {FILL2a FILL2b}} }
\end{lstlisting}

可用 \opt{-skip\_fixed\_cells}。

\subsection{用特定 Filler 邻接标准单元}
\subsubsection*{Abutting Standard Cells With Specific Filler Cells}

用 \cmd{create\_left\_right\_filler\_cells} 在特定标准单元左/右侧插入必选 filler，然后再做标准 filler 插入填剩余空隙。

\opt{-lib\_cells} 规则格式：
\begin{lstlisting}
{standard_cells} {left_filler_cells} {right_filler_cells}
\end{lstlisting}

要求：至少一侧有 filler；filler 的 \texttt{design\_type} 为 \texttt{lib\_cell} 或 \texttt{filler}；与标准单元同 site；标准单元高度为 filler 高度整数倍；宜从大到小。通配符顺序与 \cmd{get\_lib\_cells} 一致。

\begin{lstlisting}
icc2_shell> create_left_right_filler_cells \
   -lib_cells {
      { {mylib/SC1} {mylib/LF4X mylib/LF2X} {mylib/RF2X mylib/RF1X} }
      { {mylib/SC2} {mylib/LF2X mylib/LF1X} {} } }
\end{lstlisting}

插入时检查合法位置，但不检查最小 VT、OD、TPO 等先进规则；紧贴标准单元插入，不考虑单元间距规则。

\subsection{控制基于单元的 Filler 插入}
\subsubsection*{Controlling Cell-Based Filler Cell Insertion}

\cmd{create\_left\_right\_filler\_cells} 还可控制：
\begin{itemize}
  \item 区域：\opt{-boundaries}、\opt{-voltage\_areas}；
  \item 方向：默认按行允许方向；\opt{-follow\_stdcell\_orientation} 跟随标准单元（翻转时左右互换）；
  \item 一单位 tile 间隙：\opt{-rules no\_1x}（此时勿在列表中包含宽度为 1 tile 的 filler）；
  \item 命名约定：同标准 filler（\texttt{xofiller!...}）。
\end{itemize}

\subsection{基于阈值电压的 Filler 插入}
\subsubsection*{Threshold-Voltage-Based Filler Cell Insertion}

步骤：
\begin{enumerate}
  \item 确保 VT 类型信息可用；
  \item 用 \cmd{set\_vt\_filler\_rule} 定义规则；
  \item （可选）为多种邻接 VT 组合分别设规则；
  \item 用 \cmd{create\_vtcell\_fillers} 插入。
\end{enumerate}

\begin{lstlisting}
icc2_shell> set_vt_filler_rule -vt_type {default default} \
   -filler_cells {myLib/filler2x myLib/filler1x}
icc2_shell> set_vt_filler_rule -vt_type {vtA vtA} \
   -filler_cells {myLib/fillerA2x myLib/fillerA1x}
icc2_shell> create_vtcell_fillers
\end{lstlisting}

\subsection{控制基于阈值电压的 Filler 插入}
\subsubsection*{Controlling Threshold-Voltage-Based Filler Cell Insertion}

默认填充整个 block 水平标准单元行空隙。可控制 \opt{-region}、\opt{-voltage\_area}，以及命名（默认 \texttt{xofiller!...}；\opt{-prefix}/\opt{-separator}）。

\subsection{移除阈值电压 Filler 信息}
\subsubsection*{Removing the Threshold-Voltage Filler Cell Information}

\begin{lstlisting}
icc2_shell> create_vtcell_fillers -clear_vt_information
\end{lstlisting}

###{移除 Filler Cell}
\subsubsection*{Removing Filler Cells}

\begin{lstlisting}
icc2_shell> remove_cells [get_cells xofiller*]
\end{lstlisting}

仅移除有 DRC 违例的标准流程 filler：用 \cmd{remove\_stdcell\_fillers\_with\_violation}。

% ============================================================
\section{插入 Metal Fill}
\subsection*{Inserting Metal Fill}

布线后，可用金属导线填充 block 空隙以满足金属密度规则。插入前 block 应接近满足时序，且 DRC 很少或没有。

运行 \cmd{signoff\_create\_metal\_fill}，详见第~\ref{chap:icv}~章「使用 IC Validator In-Design 插入 Metal Fill」。

\begin{noteBox}
运行 \cmd{signoff\_create\_metal\_fill} 需要 IC Validator license。
\end{noteBox}

% 章末（原书约 p.544；第 7 章 IC Validator In-Design 之前）
